% Generated draft: requires mathematical and source review.
\documentclass[11pt]{article}
\usepackage[margin=25mm]{geometry}
\usepackage{fontspec}
\setmainfont{Latin Modern Roman}
\usepackage{amsmath,amssymb,mathtools,unicode-math}
\setmathfont{Latin Modern Math}
\usepackage{enumitem,needspace,xurl,hyperref}
\setlength{\parindent}{0pt}
\setlength{\parskip}{5pt}
\setlength{\emergencystretch}{3em}
\newcommand{\N}{\mathbb{N}}
\newcommand{\Z}{\mathbb{Z}}
\newcommand{\Q}{\mathbb{Q}}
\newcommand{\R}{\mathbb{R}}
\newcommand{\C}{\mathbb{C}}
\newcommand{\F}{\mathbb{F}}
\newcommand{\A}{\mathbb{A}}
\newcommand{\PP}{\mathbb{P}}
\newcommand{\E}{\mathbb{E}}
\newcommand{\eps}{\varepsilon}
\newcommand{\dd}{\,\mathrm{d}}
\begin{document}
\section*{UnsolvedMath ID: 11100065}
\textbf{Source problem code:} catalogue entry
\subsection{Definitions and mathematical statement}
Let \(E\) be an \(A_\infty\) ring spectrum. Choose the intended space \(\operatorname{Aut}(E)\) of self-equivalences with its composition structure, and form a homotopy-coherent classifying space \(B\operatorname{Aut}(E)\). The problem, attributed to Mike Hopkins, is to construct a generalized Thom spectrum from each map
\[
 f:X\longrightarrow B\operatorname{Aut}(E),
\]
and give the precise twisted half-smash interpretation suggested by the source.

The meaning of self-equivalence must be fixed. Equivalences of \(A_\infty\) ring spectra, equivalences of underlying spectra and equivalences of rank-one \(E\)-modules give different automorphism spaces. The unit-space version uses \(\operatorname{GL}_1(E)\), the invertible components of \(\Omega^\infty E\), interpreted as module self-equivalences with the chosen side convention; it is not the same as the ring-automorphism space.

A precise local-system formulation would associate to \(f\) an \(E\)-fiber local system \(\mathcal E_f\) and define its Thom object by a homotopy colimit over \(X\), in spectra or \(E\)-modules according to the automorphism choice. For a trivial local system the expected normalization is \(E\wedge\Sigma^\infty_+X\); a nontrivial one is a twisted construction, not automatically that ordinary smash. Specify functoriality, coherence, coefficient/module target and comparison to classical Thom spectra. The source explicitly says Hopkins's more detailed description was forgotten, so these choices are not supplied by the excerpt.

The summary reports a rigorous units-based construction in parametrized/infinity-categorical work. That report does not resolve which of the above literal self-equivalence meanings the original wording intended.
\subsection{Short English statement}
Construct Thom spectra for coherent bundles of copies of an \(A_\infty\) ring spectrum, clarifying the automorphism space and the twisted half-smash construction, especially the units/module version.
\subsection{Sources}
\begin{itemize}
\item Mark Hovey, \emph{Algebraic Topology Problem List}, ``Miscellaneous problems,'' item 3, explicitly attributed to Mike Hopkins, surviving mirror \url{https://www-users.cse.umn.edu/~tlawson/hovey/}. Recorded access 2026-07-29; extraction html-ordered-list.
\item The supplied summary credits Ando--Blumberg--Gepner--Hopkins--Rezk with parametrized/infinity-categorical realization, but gives no exact paper title or bibliography entry.
\item Local UnsolvedMath registry, \texttt{data/unsolvedmath/problems.json}, permanent ID \(11100065\), exact source code \texttt{catalogue entry}; supplied solved/SOLVED-IN-LITERATURE status.
\end{itemize}
\end{document}
